\documentclass[12pt]{article}
\usepackage{amsmath, amssymb, amsthm, geometry}
\geometry{a4paper, margin=2.5cm}

% Настройки для русского языка (XeLaTeX)
\usepackage{polyglossia}
\setdefaultlanguage{russian}
\setotherlanguage{english}
\usepackage{fontspec}
\setmainfont{Times New Roman}  % при наличии можно заменить на CMU Serif / Liberation Serif

\newtheorem{theorem}{Теорема}
\newtheorem{corollary}{Следствие}
\newtheorem{lemma}{Лемма}
\newtheorem{proposition}{Предложение}
\newtheorem{definition}{Определение}
\newtheorem{remark}{Замечание}

\DeclareMathOperator{\sign}{sign}
\newcommand{\R}{\mathbb{R}}

\begin{document}

\title{О невозможности сильного понижения размерности\\ в псевдоевклидовом аналоге леммы Джонсона--Линденштраусса}
\author{}
\date{}
\maketitle

\begin{abstract}
Доказывается, что для любого фиксированного \(0<\varepsilon<1\) в общем случае не существует \(\varepsilon\)-изометрического вложения (в смысле относительной ошибки квадрата псевдоевклидова расстояния) из \(\R^{p,q}\) в \(\R^{p',q'}\) с сигнатурой \(p'+q' = o(\min(p,q))\). Иными словами, любое такое вложение требует целевой сигнатуры, линейной по \(\min(p,q)\). Ключевой ингредиент: существование симметричных матриц нулей и единиц («матриц светового конуса»), не реализуемых точками пространств малой сигнатуры; нижняя оценка получается подсчётом знаковых условий квадратичных полиномов (теорема Милнора--Тома в форме Бэзу--Поллака--Роя), а точная верхняя оценка \((n,n)\) - явной грам-матричной конструкцией.
\end{abstract}

\section{Постановка задачи}

Пусть \(\R^{p,q}\) - псевдоевклидово пространство с квадратичной формой
\[
\|x\|^2 = \sum_{i=1}^p x_i^2 - \sum_{j=p+1}^{p+q} x_j^2 .
\]
Величина \(\|x-y\|^2\) может быть положительной, нулевой или отрицательной; множество пар с \(\|x-y\|^2=0\) образует \emph{световой конус}.

\begin{definition}
Пусть \(A \in \{0,1\}^{n\times n}\) - симметричная матрица с нулями на диагонали. Будем говорить, что \(A\) \emph{реализуема} в \(\R^{p,q}\), если существуют точки \(x_1,\dots,x_n \in \R^{p,q}\) такие, что
\[
A_{ij}=0 \iff \|x_i - x_j\|^2 = 0 \qquad (i\ne j).
\]
(Нули в \(A\) соответствуют нулевым, единицы - ненулевым квадратам псевдорасстояний.)
\end{definition}

Наша цель: доказать, что существует абсолютная константа \(C>0\) такая, что при \(p+q < C n\) не все матрицы \(A\) реализуемы (Теорема~\ref{thm:counting}). Затем этот факт применяется к вопросу о существовании \(\varepsilon\)-изометрических вложений (Теорема~\ref{thm:jl}).

\section{Реализуемость любой матрицы в сигнатуре \texorpdfstring{$(n,n)$}{(n,n)}}

\begin{proposition}\label{prop:realize}
Любая симметричная матрица \(A \in \{0,1\}^{n\times n}\) с нулями на диагонали реализуема в \(\R^{n,n}\).
\end{proposition}

\begin{proof}
Пусть \(B\) - симметричная матрица «дополнения»: \(B_{ij}=1\), если \(A_{ij}=1\) (т.е.\ пара должна иметь ненулевой квадрат расстояния), и \(B_{ij}=0\) иначе; \(B_{ii}=0\).

Напомним, что \emph{спектральным радиусом} симметричной матрицы \(B\) называется величина
\[
\rho(B) = \max\{\,|\lambda| : \lambda \text{ - собственное значение } B\,\}.
\]
Поскольку \(B\) вещественна и симметрична, все её собственные значения вещественны, и \(\rho(B)\) корректно определён. Выберем
\[
0<\varepsilon_0<\frac{1}{\rho(B)+1}
\]
(в частности, \(\varepsilon_0\,\rho(B)<1\); знаменатель положителен, так что такой \(\varepsilon_0\) существует). Рассмотрим матрицу
\[
M = I_n + \varepsilon_0 B .
\]

\textbf{Собственные значения \(M\).} Пусть \(v\ne 0\) - собственный вектор \(B\) с собственным значением \(\lambda\), т.е.\ \(Bv=\lambda v\). Тогда
\[
M v = (I_n+\varepsilon_0 B)v = v + \varepsilon_0\lambda v = (1+\varepsilon_0\lambda)\,v,
\]
то есть тот же вектор \(v\) является собственным для \(M\) с собственным значением \(1+\varepsilon_0\lambda\). Так как \(B\) симметрична, у неё есть ортонормированный базис из собственных векторов, и указанное соответствие \(\lambda\mapsto 1+\varepsilon_0\lambda\) исчерпывает все собственные значения \(M\).

\textbf{Положительная определённость \(M\).} Для каждого собственного значения \(\lambda\) матрицы \(B\) имеем \(|\lambda|\le\rho(B)\), поэтому \(\varepsilon_0\lambda\ge-\varepsilon_0\rho(B)>-1\), откуда \(1+\varepsilon_0\lambda>0\). Значит, все собственные значения \(M\) положительны; будучи симметричной с положительными собственными значениями, \(M\) положительно определена, \(M\succ0\).

\textbf{\(M\) - матрица Грама.} Всякая симметричная положительно определённая матрица \(M\in\R^{n\times n}\) является матрицей Грама некоторого набора векторов. Действительно, по разложению Холецкого (или через \(M=M^{1/2}M^{1/2}\)) существует матрица \(L\in\R^{n\times n}\) с \(M=L^{\top}L\); взяв в качестве \(y_1,\dots,y_n\in\R^n\) столбцы \(L\), получаем \(\langle y_i,y_j\rangle = (L^{\top}L)_{ij} = M_{ij}\).

\textbf{Свойства \(y_i\).} Из \(\langle y_i,y_j\rangle=M_{ij}\) и определения \(M=I_n+\varepsilon_0 B\) сразу следует:
\[
\|y_i\|^2=\langle y_i,y_i\rangle = M_{ii}=1,\qquad
\langle y_i,y_j\rangle = M_{ij} = \varepsilon_0 B_{ij} =
\begin{cases}
0, & A_{ij}=0\ (B_{ij}=0),\\
\varepsilon_0\ne 0, & A_{ij}=1\ (B_{ij}=1)
\end{cases}
\quad (i\ne j),
\]
где использовано \(B_{ii}=0\) (диагональ) и \(0<\varepsilon_0\).

Возьмём стандартный ортонормированный базис \(e_1,\dots,e_n\) в положительной части \(\R^n\) и положим
\[
x_i = (e_i,\; y_i) \in \R^{n}\oplus\R^{n} = \R^{n,n},\qquad
\|(u,v)\|^2 = \|u\|^2-\|v\|^2 .
\]
Тогда при \(i\ne j\)
\[
\|x_i - x_j\|^2 = \|e_i-e_j\|^2 - \|y_i-y_j\|^2
= 2 - \bigl(2 - 2\langle y_i,y_j\rangle\bigr) = 2\langle y_i,y_j\rangle .
\]
Значит, \(\|x_i-x_j\|^2 = 0\) при \(A_{ij}=0\) и \(\|x_i-x_j\|^2 = 2\varepsilon_0 \ne 0\) при \(A_{ij}=1\). Таким образом, \(A_{ij}=0\iff\|x_i-x_j\|^2=0\), и \(A\) реализована в \(\R^{n,n}\).
\end{proof}

\begin{remark}
Конструкция самодостаточна и не опирается на внешние теоремы о вложении графов. Сигнатура \((n,n)\) с точностью до постоянного множителя оптимальна: по Теореме~\ref{thm:counting} никакая сигнатура порядка \(o(n)\) не годится для всех \(A\).
\end{remark}

\section{Существование «плохих» матриц для малых размерностей}

Покажем, что не все матрицы реализуемы в пространствах с \(p+q < C n\).

Пусть \(s = p+q\) и \(D = n s\) - число вещественных координат всех точек \(x_1,\dots,x_n\in\R^{p,q}\). Для каждой пары \(i<j\) рассмотрим квадратичный полином
\[
f_{ij}(X) = \|x_i - x_j\|^2,\qquad X\in\R^{D},
\]
степени \(d=2\); всего таких полиномов \(m = \binom n2\).

Для данной матрицы \(A\) её реализуемость равносильна существованию точки \(X\in\R^D\), в которой \(f_{ij}(X)=0\) при \(A_{ij}=0\) и \(f_{ij}(X)\ne 0\) при \(A_{ij}=1\). Поскольку значения \(f_{ij}\) в псевдоевклидовом случае могут быть и отрицательными, каждой реализуемой матрице отвечает \emph{хотя бы одно} непустое знаковое условие семейства \(\{f_{ij}\}\) со значениями в \(\{-1,0,+1\}\) (единицы \(A\) соответствуют знакам \(\pm1\)). Поэтому
\[
\#\{\text{реализуемых матриц } A\}\ \le\ \#\{\text{непустых знаковых условий семейства } \{f_{ij}\}\}.
\]
Число непустых знаковых условий \(m\) полиномов степени \(\le d\) от \(D\) переменных ограничено сверху (теорема Милнора--Тома / оценка Уоррена в форме \cite[гл.~7]{BasuPollackRoy}):
\[
\#\{\text{знаковых условий}\}\ \le\ \binom{m}{\le D}\,(O(d))^{D}\ \le\ \left(\frac{C_0\, d\, m}{D}\right)^{D}
\]
для некоторой абсолютной константы \(C_0>0\) при условии \(m \ge D\).

Проверим условие \(m\ge D\): оно означает \(\binom n2 \ge ns\), т.е.\ \(s \le \tfrac{n-1}{2}\); ниже мы работаем с \(s=O(n)\) при малом коэффициенте, так что это выполнено. Подставляя \(d=2\), \(m=\tfrac{n(n-1)}2\), \(D=ns\), получаем
\[
\frac{C_0\, d\, m}{D} = \frac{C_0\,(n-1)}{s} \le \frac{C_0\, n}{s},
\qquad\text{откуда}\qquad
\#\{\text{реализуемых } A\} \le \left(\frac{C_0 n}{s}\right)^{n s}.
\]

Общее число матриц равно \(2^{\binom n2}=2^{n(n-1)/2}\). Если бы все матрицы были реализуемы, то
\[
\left(\frac{C_0 n}{s}\right)^{n s} \ge 2^{n(n-1)/2}.
\]
Логарифмируя по основанию \(2\) и деля на \(n\), имеем
\[
s\,\log_2\!\left(\frac{C_0 n}{s}\right) \ge \frac{n-1}{2}. \tag{1}
\]

Выберем \(\alpha>0\) настолько малым, чтобы одновременно
\[
\alpha\,\log_2\!\left(\frac{C_0}{\alpha}\right) < \frac14
\qquad\text{и}\qquad
\alpha < \frac{C_0}{e}. \tag{2}
\]
Положим \(g(s) = s\log_2(C_0 n/s)\). Так как \(g'(s)=\log_2(C_0 n/s)-\tfrac1{\ln 2}>0\) при \(s<C_0 n/e\), функция \(g\) возрастает на \((0,\alpha n]\) в силу второго условия~(2). Поэтому для всех \(s\le \alpha n\)
\[
g(s)\ \le\ g(\alpha n)\ =\ \alpha n\,\log_2\!\left(\frac{C_0}{\alpha}\right)\ <\ \frac n4\ \le\ \frac{n-1}{2}\quad(n\ge 2).
\]
Значит, при \(s\le\alpha n\) (и \(n\ge 3\)) неравенство~(1) \emph{не выполняется}, и найдётся матрица \(A\), не реализуемая в \(\R^{p,q}\). Полагая \(C=\alpha\), получаем следующее.

\begin{theorem}\label{thm:counting}
Существует абсолютная константа \(C>0\) такая, что для любых достаточно больших \(n\) и любых \(p,q\) с \(p+q < C n\) найдётся матрица \(A\in\{0,1\}^{n\times n}\), не реализуемая в \(\R^{p,q}\).
\end{theorem}

\section{Применение к лемме Джонсона--Линденштраусса}

Естественным псевдоевклидовым аналогом \(\varepsilon\)-изометрии служит сохранение квадрата расстояния с \emph{относительной} ошибкой; в отличие от евклидова случая, форма знакопеременна, поэтому оценка берётся по модулю.

\begin{definition}
Отображение \(f:\R^{p,q}\to\R^{p',q'}\) называется \(\varepsilon\)-изометрией (\(0<\varepsilon<1\)), если для всех \(x,y\)
\[
\bigl|\,\|f(x)-f(y)\|^2 - \|x-y\|^2\,\bigr| \le \varepsilon\,\bigl|\|x-y\|^2\bigr|. \tag{3}
\]
\end{definition}

Условие~(3) корректно при любом знаке \(\|x-y\|^2\). Для \(\|x-y\|^2\ge 0\) оно равносильно
\((1-\varepsilon)\|x-y\|^2\le \|f(x)-f(y)\|^2\le (1+\varepsilon)\|x-y\|^2\),
а для \(\|x-y\|^2<0\) - двойному неравенству с переставленными границами
\((1+\varepsilon)\|x-y\|^2\le \|f(x)-f(y)\|^2\le (1-\varepsilon)\|x-y\|^2\).

\begin{lemma}\label{lem:cone}
Всякая \(\varepsilon\)-изометрия сохраняет световой конус: для любых \(x,y\)
\[
\|f(x)-f(y)\|^2 = 0 \iff \|x-y\|^2 = 0 .
\]
\end{lemma}

\begin{proof}
Если \(\|x-y\|^2=0\), то правая часть~(3) равна \(0\), поэтому \(\|f(x)-f(y)\|^2=0\).
Если же \(\|x-y\|^2\ne0\), то из~(3)
\[
\bigl|\|f(x)-f(y)\|^2\bigr| \ge \bigl|\|x-y\|^2\bigr| - \varepsilon\bigl|\|x-y\|^2\bigr| = (1-\varepsilon)\bigl|\|x-y\|^2\bigr| > 0,
\]
поскольку \(0<\varepsilon<1\). Таким образом, \(\|f(x)-f(y)\|^2\ne0\). Оба перехода доказаны и не зависят от знака \(\|x-y\|^2\).
\end{proof}

\begin{theorem}[Нижняя оценка для псевдоевклидова JL-вложения]\label{thm:jl}
Пусть \(C>0\) - константа из Теоремы~\ref{thm:counting}. Тогда для любого фиксированного \(0<\varepsilon<1\), любых достаточно больших \(p,q\) и любой \(\varepsilon\)-изометрии \(f:\R^{p,q}\to\R^{p',q'}\) выполнено
\[
p' + q' \ \ge\ C\cdot \min(p,q).
\]
\end{theorem}

\begin{proof}
Пусть \(n=\min(p,q)\); именно \(n\) играет роль числа точек \(N\) из классической леммы, поскольку строящаяся ниже жёсткая конфигурация состоит ровно из \(n\) точек. Предположим противное: \(p'+q' < C n\). По Теореме~\ref{thm:counting} существует матрица \(A\in\{0,1\}^{n\times n}\), \emph{не реализуемая} в целевом пространстве \(\R^{p',q'}\) (мы применяем теорему именно к сигнатуре \((p',q')\) отображения \(f\)).

С другой стороны, по Предложению~\ref{prop:realize} матрица \(A\) реализуема в \(\R^{n,n}\); поскольку \(p\ge n\) и \(q\ge n\), эту конструкцию можно вложить в \(\R^{p,q}\) добавлением нулевых координат. Пусть \(x_1,\dots,x_n\in\R^{p,q}\) - соответствующие точки, так что \(A_{ij}=0\iff\|x_i-x_j\|^2=0\).

Рассмотрим их образы \(f(x_1),\dots,f(x_n)\in\R^{p',q'}\). По Лемме~\ref{lem:cone}
\[
\|f(x_i)-f(x_j)\|^2 = 0 \iff \|x_i-x_j\|^2 = 0 \iff A_{ij}=0,
\]
то есть точки \(f(x_i)\) реализуют \(A\) в \(\R^{p',q'}\). Это противоречит выбору \(A\). Следовательно, \(p'+q'\ge C n = C\min(p,q)\).
\end{proof}

\begin{corollary}[Формулировка в терминах числа точек]\label{cor:points}
Обозначим через \(N=\min(p,q)\) число точек в жёсткой конфигурации (роль \(N\) из классической леммы Джонсона--Линденштраусса). Тогда любая \(\varepsilon\)-изометрия \(f:\R^{p,q}\to\R^{p',q'}\) требует целевой сигнатуры
\[
p'+q' \ge C\cdot N,
\]
\emph{линейной по числу точек} \(N\). Это прямая переформулировка Теоремы~\ref{thm:jl}: используемая конфигурация состоит ровно из \(N=\min(p,q)\) точек, и её структуру светового конуса нельзя воспроизвести в сигнатуре, меньшей \(C\,N\).
\end{corollary}

\begin{remark}
Оценка Теоремы~\ref{thm:jl} справедлива для \emph{каждого} фиксированного \(\varepsilon\in(0,1)\): именно \(\varepsilon\)-изометричность (условие~(3)) гарантирует сохранение светового конуса (Лемма~\ref{lem:cone}), без чего перенос реализации невозможен. Таким образом, для конфигурации из \(N=\min(p,q)\) точек любая \(\varepsilon\)-изометрия требует сигнатуры
\[
p'+q' = \Omega(N),
\]
то есть \emph{линейной по числу точек} \(N\). Это в резком контрасте с евклидовой леммой Джонсона--Линденштраусса \cite{JL1984}, где для фиксированного \(\varepsilon\) набор из тех же \(N\) точек вкладывается в размерность \(O\!\bigl(\varepsilon^{-2}\log N\bigr)\), лишь \emph{логарифмическую} по \(N\) и не зависящую от объемлющей размерности.
\end{remark}

\begin{thebibliography}{9}
\bibitem{BasuPollackRoy}
S. Basu, R. Pollack, M.-F. Roy, \emph{Algorithms in Real Algebraic Geometry}, 2nd ed., Springer, 2006.

\bibitem{Milnor1964}
J. Milnor, \emph{On the Betti numbers of real varieties}, Proc. Amer. Math. Soc. \textbf{15} (1964), 275--280.

\bibitem{Thom1965}
R. Thom, \emph{Sur l'homologie des vari\'et\'es alg\'ebriques r\'eelles}, in Differential and Combinatorial Topology, Princeton Univ. Press, 1965, 255--265.

\bibitem{Warren1968}
H. E. Warren, \emph{Lower bounds for approximation by nonlinear manifolds}, Trans. Amer. Math. Soc. \textbf{133} (1968), 167--178.

\bibitem{JL1984}
W. B. Johnson, J. Lindenstrauss, \emph{Extensions of Lipschitz mappings into a Hilbert space}, Contemp. Math. \textbf{26} (1984), 189--206.
\end{thebibliography}

\end{document}
